% =========================================================
% Proof stubs — Stern–Brocot Limit Theory
% Four files; each would live at:
%   volume-ii/rationals/proofs/notes/<filename>.tex
%
% File A: sb-denom-growth.tex
% File B: sb-length-zero.tex
% File C: sb-unique-point.tex
% File D: sb-path-limit.tex
% File E: sb-limit-irrational.tex
% =========================================================


% =========================================================
% FILE A: sb-denom-growth.tex
% Proof: Lemma — Denominator growth on infinite SB paths
% Source: volume-ii/rationals/notes/<topic>.tex
% =========================================================

\subsection*{Lemma --- Denominator growth on infinite Stern--Brocot paths}
\label{prf:sb-denom-growth}

\begin{remark*}[Return]
$\leftarrow$ Back to Lemma~\ref{lem:sb-denom-growth} (Denominator growth) in Notes.
\end{remark*}

\begin{proof}
\Claim Along any infinite path in the Stern--Brocot tree, the denominators
of the successive endpoint fractions form a strictly increasing sequence
tending to $\infty$.

\Given A path in the Stern--Brocot tree.  At each level $n$ the path
occupies an interval $[a_n/b_n,\; c_n/d_n]$.  By the Farey determinant
condition (Farey Neighbor Theorem, Theorem~10.61), the two endpoints satisfy
$b_n c_n - a_n d_n = 1$ at every level.

\Goal To show $b_{n+1} > b_n$ (and symmetrically $d_{n+1} > d_n$) for
every $n$, and hence $b_n \to \infty$.

\Strategy We proceed by induction on the level $n$ of the path.

\textbf{Base case} ($n = 0$). The tree is initialised from $0/1$ and $1/0$
(the sentinel fractions). Their denominators are $1$ and $0$. When the
first mediant $1/1$ is inserted, the new interval is either $[0/1, 1/1]$
with denominators $1, 1$, or $[1/1, 1/0]$ with denominators $1, 0$.
In both cases the denominator of the chosen boundary fraction is $b_0 = 1$,
which is positive. \checkbox

\textbf{Inductive step.} Assume the interval at level $n$ is
$[a_n/b_n,\; c_n/d_n]$ with $b_n, d_n \ge 1$.

\ByDef the mediant-insertion rule, the new fraction inserted at level
$n+1$ is
\[
  \frac{a_n + c_n}{b_n + d_n}.
\]
One endpoint of the new interval is this mediant, which has denominator
$b_n + d_n$.

\ByAss $b_n \ge 1$ and $d_n \ge 1$, we have
\[
  b_n + d_n > b_n \quad \text{and} \quad b_n + d_n > d_n.
\]
\Therefore the denominator of the new endpoint strictly exceeds both
$b_n$ and $d_n$. The other endpoint retains the denominator it had
at level $n$.  In either case the minimum denominator on the boundary
strictly increases. \checkbox

\Hence by induction every denominator encountered along the path is
strictly larger than its predecessor.  Since the denominators are positive
integers, a strictly increasing sequence of positive integers is unbounded.
\Therefore $b_n \to \infty$. \AsReq
\end{proof}

\begin{remark*}[Proof shape]
Induction on depth, with the key computation being $b_n + d_n > b_n$.
The unboundedness conclusion is immediate once strict monotonicity of a
positive-integer sequence is in hand.
\end{remark*}

\begin{remark*}[Dependencies]
Farey Neighbor Theorem (Theorem~10.61) --- guarantees the determinant
condition at each level, confirming the endpoints are genuine Farey
neighbors. Positivity of natural numbers (Corollary~10.26) --- ensures
denominators are at least $1$ after the first step, making the strict
inequality $b_n + d_n > b_n$ valid.
\end{remark*}


% =========================================================
% FILE B: sb-length-zero.tex
% Proof: Lemma — SB interval lengths shrink to zero
% Source: volume-ii/rationals/notes/<topic>.tex
% =========================================================

\subsection*{Lemma --- Stern--Brocot interval lengths shrink to zero}
\label{prf:sb-length-zero}

\begin{remark*}[Return]
$\leftarrow$ Back to Lemma~\ref{lem:sb-length-zero} (SB lengths to zero) in Notes.
\end{remark*}

\begin{proof}
\Claim Let $(I_n)$ be the nested intervals from the Stern--Brocot
refinement process, with $I_n = [a_n/b_n,\; c_n/d_n]$ and
$\ell_n = c_n/d_n - a_n/b_n$. Then $\ell_n \to 0$.

\Given At each level $n$, the two endpoints $a_n/b_n$ and $c_n/d_n$
are Farey neighbors (the determinant condition is preserved by mediant
insertion, as verified in the proof of the Farey Neighbor Theorem).

\Goal To show $\forall \varepsilon > 0\ \exists N\ \forall n \ge N,\ \ell_n < \varepsilon$.

\Strategy Apply the Gap Between Farey Neighbors identity, then use the
denominator growth lemma.

\ByThm{Gap Between Farey Neighbors} (Theorem~10.65), since $a_n/b_n$
and $c_n/d_n$ are Farey neighbors at every level,
\[
  \ell_n = \frac{c_n}{d_n} - \frac{a_n}{b_n} = \frac{1}{b_n d_n}.
\]

\ByThm{Denominator growth on infinite Stern--Brocot paths}
(Lemma~\ref{lem:sb-denom-growth}), $b_n \to \infty$ and $d_n \to \infty$.

\Let $\varepsilon > 0$ be arbitrary.
\ByThm{Archimedean property of $\mathbb{Q}$} (Theorem~10.42), there
exists $N \in \mathbb{N}$ such that $b_N d_N > 1/\varepsilon$.

For all $n \ge N$, the strict monotonicity of denominators gives
$b_n \ge b_N$ and $d_n \ge d_N$, so
\[
  \ell_n = \frac{1}{b_n d_n} \le \frac{1}{b_N d_N} < \varepsilon.
\]

\Hence $\ell_n \to 0$. \AsReq
\end{proof}

\begin{remark*}[Proof shape]
A direct application of two previously established results: the exact
gap formula and the denominator growth lemma.  The Archimedean property
converts the asymptotic statement into a concrete $N$.
\end{remark*}

\begin{remark*}[Dependencies]
Gap Between Farey Neighbors (Theorem~10.65) --- provides the identity
$\ell_n = 1/(b_n d_n)$. Denominator growth on infinite Stern--Brocot
paths (Lemma~\ref{lem:sb-denom-growth}) --- supplies $b_n, d_n \to \infty$.
Archimedean property of $\mathbb{Q}$ (Theorem~10.42) --- used to find
the index $N$.
\end{remark*}


% =========================================================
% FILE C: sb-unique-point.tex
% Proof: Proposition — Nested SB intervals contain unique point
% Source: volume-ii/rationals/notes/<topic>.tex
% =========================================================

\subsection*{Proposition --- Nested Stern--Brocot intervals contain a unique limit point}
\label{prf:sb-unique-point}

\begin{remark*}[Return]
$\leftarrow$ Back to Proposition~\ref{prop:sb-unique-point} (Unique limit point) in Notes.
\end{remark*}

\begin{proof}
\Claim Let $(I_n)$ with $I_n = [a_n/b_n,\; c_n/d_n]$ be the nested
closed rational intervals from the Stern--Brocot refinement process.
Then $\bigcap_{n=0}^{\infty} I_n$ contains exactly one point.

\Given The intervals are nested: $I_0 \supseteq I_1 \supseteq I_2 \supseteq \cdots$
by construction (at each step one endpoint is replaced by the mediant,
which lies strictly inside the current interval). The lengths satisfy
$\ell_n \to 0$ by Lemma~\ref{lem:sb-length-zero}.

\Goal Existence: $\bigcap I_n \ne \varnothing$.
Uniqueness: $|\bigcap I_n| = 1$.

\Strategy Apply the Monotone Convergence Theorem in $\mathbb{R}$ to both
endpoint sequences; then use $\ell_n \to 0$ to pin them to the same limit.

\textbf{Left endpoints.} The sequence $(a_n/b_n)$ is increasing
(each left endpoint is either unchanged or replaced by a larger mediant)
and bounded above by $c_0/d_0$. \ByThm{Monotone Convergence Theorem in
$\mathbb{R}$}, the supremum $\alpha = \sup_n a_n/b_n$ exists in
$\mathbb{R}$ and $a_n/b_n \to \alpha$.

\textbf{Right endpoints.} The sequence $(c_n/d_n)$ is decreasing
and bounded below by $a_0/b_0$.  By the same theorem, the infimum
$\beta = \inf_n c_n/d_n$ exists in $\mathbb{R}$ and $c_n/d_n \to \beta$.

\Therefore $\alpha \le \beta$ (since $a_n/b_n \le c_n/d_n$ for all $n$).
But
\[
  0 \le \beta - \alpha \le \ell_n = \frac{1}{b_n d_n} \to 0,
\]
so $\alpha = \beta$.

\Let $\xi = \alpha = \beta$. Then $a_n/b_n \le \xi \le c_n/d_n$ for
all $n$, so $\xi \in I_n$ for all $n$.
\Therefore $\xi \in \bigcap I_n$, establishing existence.

\textbf{Uniqueness.} Suppose $\xi, \xi' \in \bigcap I_n$.
Then for every $n$,
\[
  |\xi - \xi'| \le \ell_n = \frac{1}{b_n d_n}.
\]
Since $\ell_n \to 0$, we have $|\xi - \xi'| = 0$, so $\xi = \xi'$.

\Hence $\bigcap I_n = \{\xi\}$. \AsReq
\end{proof}

\begin{remark*}[Proof shape]
A nested-intervals argument in $\mathbb{R}$: the left endpoints form a
bounded increasing sequence and the right endpoints form a bounded decreasing
sequence; both converge; the limit length forces the two limits to coincide.
Uniqueness is a direct consequence of the length tending to zero.
\end{remark*}

\begin{remark*}[Dependencies]
Monotone Convergence Theorem in $\mathbb{R}$ (from the bounding/sequences
chapter) --- the load-bearing tool; requires completeness of $\mathbb{R}$,
which is why the proof cannot be carried out inside $\mathbb{Q}$.
Stern--Brocot interval lengths shrink to zero (Lemma~\ref{lem:sb-length-zero})
--- supplies $\ell_n \to 0$.
\end{remark*}


% =========================================================
% FILE D: sb-path-limit.tex
% Proof: Corollary — Every infinite SB path converges
% Source: volume-ii/rationals/notes/<topic>.tex
% =========================================================

\subsection*{Corollary --- Every infinite Stern--Brocot path converges}
\label{prf:sb-path-limit}

\begin{remark*}[Return]
$\leftarrow$ Back to Corollary~\ref{cor:sb-path-limit} (SB path converges) in Notes.
\end{remark*}

\begin{proof}
\Claim Both endpoint sequences $(a_n/b_n)$ and $(c_n/d_n)$ of the
Stern--Brocot refinement process targeting $x > 0$ converge to $x$.

\Given The process constructs $I_n = [a_n/b_n, c_n/d_n]$ with $x \in I_n$
for all $n$ (by the selection rule: at each step the endpoint that does not
separate $x$ from the mediant is retained, so $x$ remains in the interval).

\Strategy Apply Proposition~\ref{prop:sb-unique-point}.

\ByThm{Nested Stern--Brocot intervals contain a unique limit point}
(Proposition~\ref{prop:sb-unique-point}), there is a unique $\xi$ with
$\xi \in \bigcap I_n$ and $a_n/b_n, c_n/d_n \to \xi$.

\ByAss $x \in I_n$ for all $n$, so $x \in \bigcap I_n$.
By the uniqueness from Proposition~\ref{prop:sb-unique-point}, $\xi = x$.

\Hence $a_n/b_n \to x$ and $c_n/d_n \to x$. \AsReq
\end{proof}

\begin{remark*}[Proof shape]
A one-step corollary: the unique point in $\bigcap I_n$ is $\xi$;
$x$ also lies in every $I_n$; therefore $\xi = x$.
\end{remark*}

\begin{remark*}[Dependencies]
Nested Stern--Brocot intervals contain a unique limit point
(Proposition~\ref{prop:sb-unique-point}).
Construction rule for the Stern--Brocot refinement process (the selection
step that keeps $x$ in the interval at every stage).
\end{remark*}


% =========================================================
% FILE E: sb-limit-irrational.tex
% Proof: Corollary — Infinite paths converge to irrationals
% Source: volume-ii/rationals/notes/<topic>.tex
% =========================================================

\subsection*{Corollary --- Infinite Stern--Brocot paths converge to irrationals}
\label{prf:sb-limit-irrational}

\begin{remark*}[Return]
$\leftarrow$ Back to Corollary~\ref{cor:sb-limit-irrational} (Irrational limits) in Notes.
\end{remark*}

\begin{proof}
\Claim A positive real number $x$ is rational if and only if it
corresponds to a finite path in the Stern--Brocot tree.

\Given The Stern--Brocot tree is generated from $0/1$ and $1/0$ by
repeated mediant insertion. At each step the process asks: is $x$ equal
to, less than, or greater than the current mediant?

\Strategy We prove both implications separately.

\medskip
\noindent$(\Rightarrow)$\ \textit{Rational implies finite path.}

\Given $x = p/q \in \mathbb{Q}_{>0}$ in reduced form.

The search process compares $x$ to successive mediants.  At each step
the two boundary fractions $a_n/b_n$ and $c_n/d_n$ satisfy $b_n c_n -
a_n d_n = 1$ (Farey Neighbor Theorem, Theorem~10.61) and the mediant
has denominator $b_n + d_n$.

Suppose $x$ is never equal to a mediant. Then at every level the
inequality is strict and both $b_n$ and $d_n$ are at most $q - 1$
(otherwise a fraction with denominator $\le q$ equal to $x$ would have
been produced, contradicting the assumption that $x = p/q$ is in reduced
form with denominator exactly $q$). But Lemma~\ref{lem:sb-denom-growth}
says $b_n \to \infty$, contradiction.

\Therefore the process must encounter $x$ as a mediant at some finite
level $N$, and the path terminates there. \checkbox

\medskip
\noindent$(\Leftarrow)$\ \textit{Finite path implies rational.}

\Given The path terminates at level $N$, meaning $x$ equals the mediant
$(a_N + c_N)/(b_N + d_N)$.

\ByDef a mediant is a rational number (sum of integers over sum of
integers).
\Therefore $x \in \mathbb{Q}$. \checkbox

\medskip
Taking the contrapositive of the $(\Rightarrow)$ direction:
if $x$ is irrational then the path is infinite.
Together with $(\Leftarrow)$ this gives the equivalence. \AsReq
\end{proof}

\begin{remark*}[Proof shape]
A biconditional proof. The $(\Leftarrow)$ direction is trivial (a mediant
is visibly rational). The $(\Rightarrow)$ direction is the substantive half:
it proceeds by contradiction, using denominator growth to force $x$ to
appear as a mediant within a bounded number of steps.
\end{remark*}

\begin{remark*}[Dependencies]
Farey Neighbor Theorem (Theorem~10.61) --- the determinant condition
$b_n c_n - a_n d_n = 1$ controls what denominators can appear. Denominator
growth on infinite Stern--Brocot paths (Lemma~\ref{lem:sb-denom-growth})
--- the engine of the $(\Rightarrow)$ contradiction. Definition of reduced
form for rational numbers (Definition~10.8 and Remark~10.15).
\end{remark*}